\section*{第 \ref{ch:g6:water-underground} 章——地下水}

\begin{solution}{exo:g6:water-underground:1}
\omterm{def:g6:water-underground:permeability}{透水的}：沙、\omterm{def:g3:stones-pebbles-sand:pebble}{砾石}、白垩。\omterm{def:g6:water-underground:permeability}{不透水的}：黏土、无裂隙的花岗岩、
页岩。
\end{solution}

\begin{solution}{exo:g6:water-underground:2}
(a) \omterm{def:g6:water-underground:water-table}{地下水位}；(b) \omterm{def:g6:water-underground:aquifer}{含水层}；(c) 泉；(d) \omterm{def:g6:water-underground:spring}{水井}。
\end{solution}

\begin{solution}{exo:g6:water-underground:3}
(a) 水分蒸发回到空气中；水沿地表流入溪流；
水渗入地下（\omterm{def:g6:water-underground:infiltration}{下渗}）。(b) 在停车场上：
其坚硬的路面几乎不让水渗入，因此水流走。
在草地上，土壤和草能让大部分水渗入。
\end{solution}

\begin{solution}{exo:g6:water-underground:4}
(a) 错：黏土多孔但不透水，因为其\omterm{def:g6:water-underground:porosity}{孔隙}极小。
(b) 对。(c) 错：\omterm{def:g6:water-underground:water-table}{地下水位}随\omterm{def:g4:world-physical-map:relief}{地形起伏}，但起伏更平缓，
并且随季节升降。(d) 对：约 30~\% 的\omterm{def:g2:water-on-earth:salt-fresh}{淡水}为
\omterm{def:g6:water-underground:groundwater}{地下水}，而\omterm{def:g1:landscapes-around-us:water}{湖泊}中的\omterm{def:g2:water-on-earth:salt-fresh}{淡水}仅占约四分之一个百分点。
\end{solution}

\begin{solution}{exo:g6:water-underground:5}
$10 \times 25 \div 100 = 2.5$：溢出前可倒入 2.5~升。
$10 \times 22 \div 100 = 2.2$：2.2~升水会从孔中流出。
\end{solution}

\begin{solution}{exo:g6:water-underground:6}
(a) $2 \times 1 \times 1 = 2$~立方米，即 2000~升。
(b) $2000 \times 20 \div 100 = 400$~升。(c) $2000 \times 18 \div 100 =
360$~升流出；$400 - 360 = 40$~升留在其中。
\end{solution}

\begin{solution}{exo:g6:water-underground:7}
(a) 泉的\omterm{def:g4:world-physical-map:altitude}{海拔}约为 203~m，\omterm{def:g1:landscapes-around-us:water}{河流}约为 132~m。(b) 泉
比\omterm{def:g1:landscapes-around-us:water}{河流}高约 70~m：其水沿斜坡流入河中。
(c) 在白垩山坡内部，\omterm{def:g6:water-underground:water-table}{地下水位}低于地面；
只有白垩层底部（覆于\omterm{def:g6:water-underground:permeability}{不透水的}页岩之上）
在斜坡上露出时，\omterm{def:g6:water-underground:water-table}{地下水位}才与地表相交。
\end{solution}

\begin{solution}{exo:g6:water-underground:8}
(a) 雨季过后：$401 - 336 = 65$~m；旱季过后：
$401 - 314 = 87$~m。(b) $336 - 314 = 22$~m。
\end{solution}

\begin{solution}{exo:g6:water-underground:9}
更深。\omterm{def:g6:water-underground:water-table}{地下水位}随\omterm{def:g4:world-physical-map:relief}{地形起伏}，但起伏更平缓：在
山顶处，\omterm{def:g6:water-underground:water-table}{地下水位}比井~A 下方的高，但地面升高得多，
因此水位位于地表以下更深处（在剖面图上，
深度超过 110~m）。
\end{solution}

\begin{solution}{exo:g6:water-underground:10}
黏土的\omterm{def:g6:water-underground:porosity}{孔隙}极小：水附着在颗粒上，无法
穿过黏土流向\omterm{def:g6:water-underground:spring}{水井}。黏土多孔但不透水，因此
只有少量水（500~升中的 20~升）流出。
\end{solution}

\begin{solution}{exo:g6:water-underground:11}
(a) $30.1 \div 0.26 \approx 116$：多出一百多倍。
(b) $100 - 69 = 31$：剩余约 31~\%，其中\omterm{def:g6:water-underground:groundwater}{地下水}约占 30
个百分点。(c) 它是最大的液态\omterm{def:g2:water-on-earth:salt-fresh}{淡水}储存库；冰
冻结且大多远在极地和高山。另外，\omterm{def:g6:water-underground:groundwater}{地下水}
几乎遍布我们脚下。
\end{solution}

\begin{solution}{exo:g6:water-underground:12}
两口井都从同一白垩\omterm{def:g6:water-underground:aquifer}{含水层}取水。功率强大的水泵抽水
的速度快于雨水补给白垩层的速度，因此\omterm{def:g6:water-underground:water-table}{地下水位}
整体下降，降至第二口井井底以下，该井因此干涸。农民
可以约定少抽水（尤其是夏季），共用一口井，或将
第二口井加深；最重要的是，抽取的水量不应超过
雨水的补给量。
\end{solution}

\begin{solution}{pb:g6:water-underground:1}
\textbf{1.} 农场位于白垩层之上，其下为页岩。白垩透水，
页岩不透水。

\textbf{2.} 雨水渗入白垩层，并在页岩处停止\omterm{def:g6:water-underground:infiltration}{下渗}。白垩层底部
在斜坡上出露处，\omterm{def:g6:water-underground:water-table}{地下水位}与地面相交，
水便涌出，这就是泉。

\textbf{3.} 位于地面以下 $401 - 336 = 65$~m。

\textbf{4.} 井底位于 $401 - 100 = 301$~m 处，低于\omterm{def:g6:water-underground:water-table}{地下水位}
（314~m），井仍能取到水，比水位低 $314 - 301 =
13$~m。

\textbf{5.} $100 \times 50 \times 2 = 10\,000$~立方米。

\textbf{6.} $10\,000 \times 20 \div 100 = 2000$~立方米。

\textbf{7.} $10\,000 \times 18 \div 100 = 1800$~立方米流出；
$2000 - 1800 = 200$~立方米留在其中。

\textbf{8.} $1800 \times 1000 = 1\,800\,000$~升。

\textbf{9.} 每天 $18 \times 1000 = 18\,000$~升。

\textbf{10.} 每周 $18 \times 7 = 126$~立方米。

\textbf{11.} 雨季的雨水渗入白垩层，每年
都为其补水（前提是农场抽水量不超过雨水的补给量）。

\textbf{12.} $1800 \div 18 = $~\textbf{100 天}：略多于三
个月没有任何降雨。
\end{solution}
